\begin{titlepage}
\newgeometry{margin=0.95in}
\thispagestyle{empty}

\begin{center}
{\large University of Rochester\\}
{\normalsize Warner School of Education\\[0.35in]}

{\color{PrimaryRed}\rule{\textwidth}{1.4pt}}\\[0.22in]
{\Huge\bfseries Literacy as Power:\\[0.08in] Communication Access,\\[0.08in] Neurodivergence, and Social Justice\\[0.08in] in STEM Education}\\[0.18in]
{\Large Critical Literacy Journal -- Clean V2}\\[0.22in]
{\color{PrimaryRed}\rule{0.78\textwidth}{0.7pt}}\\[0.32in]

\colorbox{SoftBlue}{%
\parbox{0.88\textwidth}{%
\centering\vspace{0.12in}
\textbf{Communication Access as Literacy, Power, and Educational Justice}\\[0.05in]
A literacy and social justice analysis of classroom communication, disciplinary language, accommodation systems, and institutional literacy.
\vspace{0.12in}
}%
}
\end{center}

\vspace{0.25in}

\begin{center}
\renewcommand{\arraystretch}{1.15}
\setlength{\tabcolsep}{4pt}
\begin{tabular}{>{\bfseries\RaggedLeft}p{1.65in} p{4.9in}}
Student & Piter Zacari Garcia Bautista \\
Course & EDU498: Literacy Learning as Social Practice \\
Assignment & Literacy as Power: Bridging Social Justice and Critical Communication \\
Focus & Communication-access injustice for neurodivergent and chronically ill students in STEM education \\
Present-Day Connection & Universal Design for Learning / CAST UDL Guidelines 3.0 \\
Submission Date & June 2026 \\
\end{tabular}
\end{center}

\vspace{0.25in}

\begin{center}
\begin{minipage}{0.95\textwidth}
\itshape
This journal uses a BIO/STEM communication-access case as an evidence site for a broader literacy and social justice argument. It examines how classroom communication, disciplinary terminology, accommodations, and institutional pathways determine whether students are understood, supported, and allowed to participate with dignity.
\end{minipage}
\end{center}

\restoregeometry
\end{titlepage}
\clearpage
\onehalfspacing

\section*{The Bridge}

\placeholderfigure{The bridge from personal evidence to systemic literacy}{Insert: real_pain_real_injustice.png}{A visual map of communication harm, institutional rigidity, and the emotional burden created by inaccessible teaching design.}

This paper examines the intersection of literacy and social justice by showing how communication systems can empower or exclude students. In my BIO/STEM communication-access case, literacy did not only mean reading and writing. It also meant decoding requirements, grading language, accommodation pathways, team expectations, scientific communication norms, and institutional procedures. When a course communicates through one narrow model of participation, students who process, remember, communicate, or recover differently can be positioned as difficult rather than under-supported.

The case matters because it raises a larger question: if an adult graduate student with technical knowledge, documentation habits, advisor support, and self-advocacy skills can still be harmed by inaccessible communication, what happens to children and younger students who cannot yet defend themselves?

\section*{Communication Access as a Social Justice Issue -- Research Findings}

\placeholderfigure{Identifying the real topic}{Insert: identifying_the_real_topic.png}{A visual map showing how communication breakdowns and institutional rigidity can shift the burden onto students instead of addressing system-level barriers.}

The issue I am focusing on is \textit{communication-access injustice for neurodivergent and chronically ill students in STEM education}. This sits directly at the intersection of literacy and social justice because school literacy is not only the ability to read and write words. It also includes the ability to access, decode, interpret, and respond to the communication systems that organize learning.

In a STEM classroom, those systems include the syllabus, assignment instructions, grading policies, verbal updates, accommodation procedures, email exchanges, evaluations, teamwork expectations, scientific terminology, documentation, and institutional support pathways. A student is not only learning bioinformatics, computer science, or data science; the student is also learning how to navigate the course itself. When these systems are clear, stable, accessible, and multimodal, literacy empowers students. When they are vague, single-mode, inconsistent, hidden, or built around one assumed kind of learner, literacy becomes gatekeeping.

\subsection*{Why This Issue Matters}

\placeholderfigure{Why this issue matters}{Insert: why_this_issue_matters.png}{A visual explanation of how inaccessible communication can turn learning barriers into personal blame.}

Communication determines access. If a student cannot reliably find, understand, remember, or clarify the language of a course, then the student cannot fully access the learning the course claims to offer. Inaccessible communication affects more than grades; it can shape health, executive functioning, mental health, peer relationships, academic confidence, trust in teachers, willingness to seek help, and future participation in STEM.

This is especially urgent in STEM because courses can confuse rigor with rigidity. A course can be intellectually demanding without being inaccessible. Rigor should mean meaningful disciplinary challenge, not ambiguity, hidden expectations, memory-dependent instructions, or emotional toughness. Access should be part of rigor because students cannot show what they know if they must first survive a communication system that was not designed for them.

\subsection*{Relevant Sources and Evidence Scope}

\placeholderfigure{Building the evidence}{Insert: building_the_evidence.png}{A visual map showing how repeated communication barriers and documented interactions create evidence of a systemic access concern.}

The research for this paper draws from four sources of evidence: a curated BIO/STEM communication-access case file, literacy and social justice course readings, teaching-placement observations, and a present-day accessibility initiative. Together, these sources connect one documented adult learning experience to broader questions of literacy, power, disability, disciplinary participation, and educational design.

The case file provides the adult evidence. The course readings provide the conceptual framework for interpreting literacy as identity, access, disciplinary participation, criticality, and power. Teaching-placement observations provide the K--12 bridge by showing how children often demonstrate understanding through action, tools, collaboration, movement, design, peer explanation, and guided practice. Universal Design for Learning, especially CAST UDL Guidelines 3.0, supports the move from documentation to redesign by treating learner variability as expected rather than exceptional \citep{cast2024udl}.

\section*{Literacy as Communication Power -- A Critical Analysis}

\placeholderfigure{Literacy as communication power}{Insert: literacy_as_communication_power.png}{A visual framing of literacy as communication access, interpretation, belonging, and institutional power.}

Literacy shapes how communication-access injustice is understood because the issue is often hidden behind neutral-sounding language. Words like \textit{rigor}, \textit{professionalism}, \textit{participation}, \textit{responsibility}, \textit{reasonable accommodation}, and \textit{student expectations} can appear fair while still placing uneven burdens on students whose bodies, brains, languages, or trauma histories do not match the assumed classroom norm.

In the BIO/STEM case, the problem was not only that some instructions were unclear. The deeper issue was that the course required students to decode which expectations were official, how grading language should be interpreted, when verbal directions carried policy weight, how illness should be translated into institutional language, and how accommodations could be used without being treated as a burden. These are not minor communication details; they are access conditions.

\subsection*{Language, Tone, and Authority}

The language of a course carries institutional power because students rarely meet course communication as equals. A professor's wording can determine what counts as responsible, late, complete, professional, respectful, or valid. When expectations are unstable, implied, or single-mode, the student must interpret the system before they can complete the assignment. For a neurodivergent student, that ambiguity can become an executive-function crisis. For a chronically ill student, it can become another demand layered on top of symptoms and disclosure fatigue. For a child, it may simply become shame, shutdown, or punishment.

The issue, then, is not whether communication exists. The issue is whether communication is accessible, stable, interpretable, and humane.
